\documentclass[11pt,aspectratio=169]{ctexbeamer}

\usetheme{default}
\setbeamertemplate{navigation symbols}{}

\usepackage{amsmath}
\usepackage{booktabs}

\definecolor{DeepBlue}{RGB}{25,74,122}
\definecolor{SoftInk}{RGB}{53,59,66}
\definecolor{PaperBg}{RGB}{249,248,245}
\definecolor{PanelBlue}{RGB}{236,243,249}

\setbeamercolor{normal text}{fg=SoftInk,bg=PaperBg}
\setbeamercolor{structure}{fg=DeepBlue}
\setbeamercolor{frametitle}{fg=DeepBlue,bg=PaperBg}
\setbeamercolor{title}{fg=DeepBlue}
\setbeamercolor{block title}{fg=DeepBlue,bg=PanelBlue}
\setbeamercolor{block body}{fg=SoftInk,bg=white}

\newcommand{\Core}[1]{
  \begin{block}{核心结论}
    #1
  \end{block}
}

\title{大模型推理系统与实践}
\subtitle{研究生课程导论}
\author{张书豪}
\institute{华中科技大学计算机学院}
\date{}

\begin{document}

\begin{frame}[plain]
  \titlepage
\end{frame}

\begin{frame}{课程定位}
  \Core{本课程面向研究生系统讲授大模型推理基础设施的关键问题、核心机制与研究方法。}
  \begin{itemize}
    \item 关注在线服务系统，而非训练过程。
    \item 关注调度、状态管理、执行路径与异构适配。
    \item 结合论文精读、系统案例、实验训练与课程项目。
  \end{itemize}
\end{frame}

\begin{frame}{为什么值得独立成课}
  \Core{大模型推理已经从“把模型跑起来”演进为“把服务稳定运行起来”的系统问题。}
  \begin{itemize}
    \item 长上下文、MoE、RAG、智能体工作流重塑了工作负载。
    \item TTFT、P99、goodput、成本与稳定性成为核心目标。
    \item 真实收益取决于全链路而非局部算子。
  \end{itemize}
\end{frame}

\begin{frame}{统一视角：状态化系统}
  \Core{推理系统的主要瓶颈，往往来自共享状态的访问、保留、搬运和复用。}
  \begin{itemize}
    \item 模型参数之外，还有 KV Cache、知识缓存、会话记忆与调度元数据。
    \item 状态跨时间、跨请求、跨设备持续传播。
    \item 调度、缓存与执行路径因此连续耦合。
  \end{itemize}
\end{frame}

\begin{frame}{课程主线}
  \Core{课程围绕三条主线展开：请求组织与调度、状态管理与复用、执行路径优化与验证。}
  \begin{enumerate}
    \item 工作负载与评价指标
    \item 推理执行路径与系统架构
    \item 调度机制与 KV 管理
    \item 异构平台适配与论文精读
    \item 开源系统实践与课程项目
  \end{enumerate}
\end{frame}

\begin{frame}{课程模块}
  \begin{itemize}
    \item 模块 1：导论与问题地图
    \item 模块 2：工作负载与评价指标
    \item 模块 3：一次请求的完整链路
    \item 模块 4：请求调度与服务编排
    \item 模块 5：KV Cache 与状态管理
    \item 模块 6：推理系统架构
    \item 模块 7：执行优化与异构路径
    \item 模块 8：论文精读与研究训练
  \end{itemize}
\end{frame}

\begin{frame}{课程能力目标}
  \Core{课程目标不是记住名词，而是形成系统研究与实验验证能力。}
  \begin{itemize}
    \item 能分解一次推理请求的生命周期。
    \item 能分析瓶颈对象、控制点与收益链路。
    \item 能阅读系统代码并设计可信实验。
    \item 能提出具备实现路径的研究选题。
  \end{itemize}
\end{frame}

\begin{frame}{课程方法}
  \Core{机制讲解、论文精读、代码导读、实验训练和课程项目共同构成完整教学闭环。}
  \begin{itemize}
    \item 不把论文讲成摘要复述。
    \item 不把系统课讲成纯部署教程。
    \item 强调问题定义、机制边界与实验可信性。
  \end{itemize}
\end{frame}

\begin{frame}{课程总结}
  \Core{真正的大模型推理基础设施研究，是让系统在复杂负载、持续状态和异构约束下稳定地跑得对、跑得快、跑得住。}
  \begin{itemize}
    \item 这是研究问题，也是工程问题。
    \item 这也是本课程希望训练的核心能力。
  \end{itemize}
\end{frame}

\end{document}
